% META: {"id":"TSPE-DERIVATION-CONVEXITE-QCM","chapitre":"TSPE-DERIVATION-CONVEXITE","type_objet":"qcm","genere_depuis":"chapitres/TSPE-DERIVATION-CONVEXITE/qcm/TSPE-DERIVATION-CONVEXITE-QCM.json","status":"approved","origin":"generated","status_history":["generated","audited","accepted","approved"],"release_acceptance":"RELEASE_OWNER_BATCH_ACCEPTANCE","acceptance_closure_digest":"sha256:9b3ccf9a81c5520fbb7e03b7b2d3e7bf2057a02834b6d908bf3be5f49f3b553f","acceptance_receipt_digest":"sha256:4fad86f0282e0fa4e0419cca1ad6379ae7af5a280c04a4a90880f90a1c044242"}
% Fichier genere par scripts/build_qcm_tex.py — ne pas editer a la main.

\section*{\textcolor{chapcolor}{\MakeUppercase{QCM d'auto-évaluation -- Complements sur la dérivation, convexité}}}

\begin{center}
\textit{Pour chaque question, une seule réponse est exacte.}
\end{center}

\begin{enumerate}

\item[] \textbf{Capacité C1}

\item \textbf{[Q1]} La dérivée de $f(x) = \mathrm{e}^{x^2}$ est :
  \begin{enumerate}[label=\Alph*.]
    \item $2x\,\mathrm{e}^{x^2}$
    \item $\mathrm{e}^{x^2}$
    \item $x\,\mathrm{e}^{x^2}$
    \item $2\,\mathrm{e}^{x^2}$
  \end{enumerate}

\item \textbf{[Q2]} La dérivée de $f(x) = \ln(x^2 + 1)$ sur $\mathbb{R}$ est :
  \begin{enumerate}[label=\Alph*.]
    \item $\dfrac{1}{x^2 + 1}$
    \item $\dfrac{2x}{x^2 + 1}$
    \item $\dfrac{2}{x^2 + 1}$
    \item $\dfrac{x}{x^2 + 1}$
  \end{enumerate}

\item \textbf{[Q3]} La dérivée de $f(x) = (3x - 2)^4$ est :
  \begin{enumerate}[label=\Alph*.]
    \item $4(3x-2)^3$
    \item $3(3x-2)^4$
    \item $12(3x-2)^3$
    \item $12(3x-2)^4$
  \end{enumerate}

\item[] \textbf{Capacité C2}

\item \textbf{[Q4]} Soit $f(x) = x\,\mathrm{e}^{-x}$. La dérivée $f'(x)$ est :
  \begin{enumerate}[label=\Alph*.]
    \item $\mathrm{e}^{-x} - x\,\mathrm{e}^{-x}$
    \item $\mathrm{e}^{-x} + x\,\mathrm{e}^{-x}$
    \item $-\mathrm{e}^{-x}$
    \item $x\,\mathrm{e}^{-x}$
  \end{enumerate}

\item \textbf{[Q5]} La limite de $f(x) = \dfrac{\ln x}{x}$ en $+\infty$ est :
  \begin{enumerate}[label=\Alph*.]
    \item $+\infty$
    \item $0$
    \item $1$
    \item $-\infty$
  \end{enumerate}

\item \textbf{[Q6]} La fonction $f(x) = x^3 - 3x$ admet comme extrema locaux :
  \begin{enumerate}[label=\Alph*.]
    \item un maximum local en $x = 1$ et un minimum local en $x = -1$
    \item un maximum local en $x = 0$
    \item un maximum local en $x = -1$ et un minimum local en $x = 1$
    \item aucun extremum local
  \end{enumerate}

\item[] \textbf{Capacité C5}

\item \textbf{[Q7]} La fonction $f(x) = \mathrm{e}^x$ est :
  \begin{enumerate}[label=\Alph*.]
    \item ni convexe ni concave
    \item concave sur $\mathbb{R}$
    \item convexe sur $\mathbb{R}^+$ et concave sur $\mathbb{R}^-$
    \item convexe sur $\mathbb{R}$
  \end{enumerate}

\item[] \textbf{Capacité C3}

\item \textbf{[Q8]} L'inégalité $\mathrm{e}^x \geq 1 + x$ est valable :
  \begin{enumerate}[label=\Alph*.]
    \item pour tout $x \in \mathbb{R}$
    \item pour tout $x \leq 0$
    \item pour tout $x \geq 0$
    \item seulement en $x = 0$
  \end{enumerate}

\item[] \textbf{Capacité C5}

\item \textbf{[Q9]} La fonction $f(x) = \ln x$ est :
  \begin{enumerate}[label=\Alph*.]
    \item convexe sur $]0, +\infty[$
    \item concave sur $]0, +\infty[$
    \item convexe puis concave
    \item ni convexe ni concave
  \end{enumerate}

\item[] \textbf{Capacité C4}

\item \textbf{[Q10]} Si $f''(x) > 0$ sur un intervalle $I$, alors la courbe de $f$ est :
  \begin{enumerate}[label=\Alph*.]
    \item croissante
    \item concave (en forme de $\cap$)
    \item convexe (en forme de $\cup$)
    \item décroissante
  \end{enumerate}

\item[] \textbf{Capacité C5}

\item \textbf{[Q11]} Un point d'inflexion est un point ou :
  \begin{enumerate}[label=\Alph*.]
    \item $f'(x) = 0$
    \item $f(x) = 0$
    \item $f''(x) = 0$ seulement
    \item la convexité change au passage de ce point
  \end{enumerate}

\item \textbf{[Q12]} Si $f$ est convexe, alors la courbe est :
  \begin{enumerate}[label=\Alph*.]
    \item au-dessus de ses tangentes
    \item au-dessous de ses tangentes
    \item au-dessus de ses cordes
    \item au-dessus de ses tangentes et de ses cordes
  \end{enumerate}

\item \textbf{[Q13]} La courbe de $f(x) = x^3$ traverse sa tangente en $x = 0$ car :
  \begin{enumerate}[label=\Alph*.]
    \item $f'(0) = 0$
    \item $f''(0) = 0$ et $f''$ change de signe
    \item $f(0) = 0$
    \item $f$ est décroissante
  \end{enumerate}

\item[] \textbf{Capacité C6}

\item \textbf{[Q14]} Si $f''(x) \geq 0$ sur $I$, alors pour tout $a, x \in I$ :
  \begin{enumerate}[label=\Alph*.]
    \item $f(x) = f(a) + f'(a)(x - a)$
    \item $f(x) \leq f(a) + f'(a)(x - a)$
    \item $f(x) \geq f(a) + f'(a)(x - a)$
    \item $f(x) \geq f(a)$
  \end{enumerate}

\item \textbf{[Q15]} L'approximation $\mathrm{e}^{0{,}1} \approx 1{,}1$ obtenue par tangente en $0$ est :
  \begin{enumerate}[label=\Alph*.]
    \item impossible a déterminer
    \item par exces (la vraie valeur est plus petite)
    \item exacte
    \item par defaut (la vraie valeur est plus grande)
  \end{enumerate}

\end{enumerate}

% NEXUS-QCM-TEACHER-ONLY-BEGIN
\ifnxVersionProfesseur
\clearpage
\section*{Clé de correction — réservée au professeur}

\begin{center}
\begin{tabular}{lll}
\hline
Question & Capacité & Réponse exacte \\
\hline
Q1 & C1 & \textbf{A} \\
Q2 & C1 & \textbf{B} \\
Q3 & C1 & \textbf{C} \\
Q4 & C2 & \textbf{A} \\
Q5 & C2 & \textbf{B} \\
Q6 & C2 & \textbf{C} \\
Q7 & C5 & \textbf{D} \\
Q8 & C3 & \textbf{A} \\
Q9 & C5 & \textbf{B} \\
Q10 & C4 & \textbf{C} \\
Q11 & C5 & \textbf{D} \\
Q12 & C5 & \textbf{A} \\
Q13 & C5 & \textbf{B} \\
Q14 & C6 & \textbf{C} \\
Q15 & C6 & \textbf{D} \\
\hline
\end{tabular}
\end{center}

\subsection*{Diagnostic des réponses erronées}

\begin{itemize}
  \item \textbf{Q1}
  \begin{itemize}
    \item \textbf{B} — La dérivée de l'exponentielle est conservee, mais la dérivée de la fonction interieure $x^2$ fournit aussi le facteur $2x$. \emph{Renvoi : C1.}
    \item \textbf{C} — Le facteur $x$ correspond a une dérivée incomplete de $x^2$ : on a $(x^2)'=2x$, et non simplement $x$. \emph{Renvoi : C1.}
    \item \textbf{D} — Le facteur $2$ oublie la variable dans $(x^2)'=2x$ ; la règle de chaine donne donc $2x\,\mathrm{e}^{x^2}$. \emph{Renvoi : C1.}
  \end{itemize}
  \item \textbf{Q2}
  \begin{itemize}
    \item \textbf{A} — La formule $(\ln u)'=u'/u$ exige le numérateur $u'=2x$ ; choisir $1$ revient a oublier toute la dérivée de $x^2+1$. \emph{Renvoi : M1, dérivée de $\ln(u)$.}
    \item \textbf{C} — Le numérateur $2$ remplace a tort $(x^2)'=2x$ par une constante et perd donc le facteur variable $x$. \emph{Renvoi : M1, dérivée de $\ln(u)$.}
    \item \textbf{D} — Le numérateur $x$ oublie le coefficient $2$ dans $(x^2+1)'=2x$ ; le dénominateur seul ne corrige pas cet oubli. \emph{Renvoi : M1, dérivée de $\ln(u)$.}
  \end{itemize}
  \item \textbf{Q3}
  \begin{itemize}
    \item \textbf{A} — Le facteur extérieur $4$ est present, mais la dérivée de la fonction interieure $3x-2$ apporte aussi le facteur $3$. \emph{Renvoi : C1.}
    \item \textbf{B} — Le facteur $3$ vient de la dérivée de $3x-2$, mais la puissance doit diminuer de $4$ a $3$ et être multipliée par $4$. \emph{Renvoi : C1.}
    \item \textbf{D} — Le coefficient $12$ est correct, mais la dérivation de la puissance quatrieme impose $(3x-2)^3$, et non $(3x-2)^4$. \emph{Renvoi : C1.}
  \end{itemize}
  \item \textbf{Q4}
  \begin{itemize}
    \item \textbf{B} — Le signe plus suppose a tort que $(\mathrm{e}^{-x})'=\mathrm{e}^{-x}$ ; la dérivée de $-x$ impose le terme $-x\mathrm{e}^{-x}$. \emph{Renvoi : M5, dérivée d'un produit avec exponentielle.}
    \item \textbf{C} — Le résultat $-\mathrm{e}^{-x}$ ne conserve que la dérivée de l'exponentielle et omet le terme obtenu en derivant le facteur $x$. \emph{Renvoi : M5, dérivée d'un produit avec exponentielle.}
    \item \textbf{D} — Le produit $x\mathrm{e}^{-x}$ est la fonction initiale : aucun des deux termes de la formule $(uv)'=u'v+uv'$ n'a été calcule. \emph{Renvoi : M5, dérivée d'un produit avec exponentielle.}
  \end{itemize}
  \item \textbf{Q5}
  \begin{itemize}
    \item \textbf{A} — Choisir $+\infty$ ne compare que les divergences separees de $\ln x$ et de $x$ ; le dénominateur croit plus vite et impose la limite nulle. \emph{Renvoi : M2, croissances comparées en $+\infty$.}
    \item \textbf{C} — La valeur $1$ supposerait des croissances equivalentes ; or $\ln x=o(x)$, donc leur quotient tend vers $0$. \emph{Renvoi : M2, croissances comparées en $+\infty$.}
    \item \textbf{D} — Pour $x>1$, $\ln x$ et $x$ sont positifs : leur quotient ne peut pas tendre vers $-\infty$, et la croissance comparee donne $0$. \emph{Renvoi : M2, croissances comparées en $+\infty$.}
  \end{itemize}
  \item \textbf{Q6}
  \begin{itemize}
    \item \textbf{A} — Cette reponse inverse les deux extrema : $f'=3(x^2-1)$ est positive, puis négative, puis positive, donc le maximum est en $-1$. \emph{Renvoi : M2, extrema et signe de la dérivée.}
    \item \textbf{B} — Le point $0$ n'est pas critique puisque $f'(0)=-3$ ; il ne peut donc pas être le maximum local annonce. \emph{Renvoi : M2, extrema et signe de la dérivée.}
    \item \textbf{D} — La dérivée $3(x^2-1)$ change deux fois de signe, en $-1$ puis en $1$ : la fonction admet donc bien deux extrema locaux. \emph{Renvoi : M2, extrema et signe de la dérivée.}
  \end{itemize}
  \item \textbf{Q7}
  \begin{itemize}
    \item \textbf{A} — La dérivée seconde $f''(x)=\mathrm{e}^x$ est positive partout ; ce critere etablit directement la convexité sur $\mathbb{R}$. \emph{Renvoi : M4, convexité par le signe de la dérivée seconde.}
    \item \textbf{B} — La concavité exigerait $f''(x)\leq 0$ ; ici $f''(x)=\mathrm{e}^x>0$ sur tout $\mathbb{R}$. \emph{Renvoi : M4, convexité par le signe de la dérivée seconde.}
    \item \textbf{C} — Le signe de $x$ ne change pas celui de $f''(x)=\mathrm{e}^x$, qui reste strictement positif sur les deux demi-droites. \emph{Renvoi : M4, convexité par le signe de la dérivée seconde.}
  \end{itemize}
  \item \textbf{Q8}
  \begin{itemize}
    \item \textbf{B} — Restreindre l'inégalité a $x\leq0$ oublie la même propriété de tangente, valable sur tout le domaine de l'exponentielle. \emph{Renvoi : M3, position de la courbe et de sa tangente.}
    \item \textbf{C} — Restreindre l'inégalité a $x\geq0$ oublie que la convexité de l'exponentielle place sa courbe au-dessus de la tangente en $0$ sur tout $\mathbb{R}$. \emph{Renvoi : M3, position de la courbe et de sa tangente.}
    \item \textbf{D} — L'égalité a lieu en $x=0$, mais l'inégalité reste stricte pour tout $x\neq0$ car l'exponentielle est strictement convexe. \emph{Renvoi : M3, position de la courbe et de sa tangente.}
  \end{itemize}
  \item \textbf{Q9}
  \begin{itemize}
    \item \textbf{A} — La convexité exigerait $f''(x)\geq0$ ; pour $f(x)=\ln x$, on a $f''(x)=-1/x^2<0$ sur $]0,+\infty[$. \emph{Renvoi : M4, convexité par le signe de la dérivée seconde.}
    \item \textbf{C} — Aucun changement de convexité ne se produit : $f''(x)=-1/x^2$ reste strictement négatif sur tout le domaine. \emph{Renvoi : M4, convexité par le signe de la dérivée seconde.}
    \item \textbf{D} — Le signe constant $f''(x)=-1/x^2<0$ fournit precisement une conclusion : $\ln$ est concave sur $]0,+\infty[$. \emph{Renvoi : M4, convexité par le signe de la dérivée seconde.}
  \end{itemize}
  \item \textbf{Q10}
  \begin{itemize}
    \item \textbf{A} — Le signe de $f''$ décrit la variation de $f'$ et la convexité, pas le signe de $f'$ nécessaire pour conclure que $f$ est croissante. \emph{Renvoi : M4, convexité par le signe de la dérivée seconde.}
    \item \textbf{B} — La concavité correspond au signe oppose $f''\leq0$ ; une dérivée seconde strictement positive caractérise ici la convexité. \emph{Renvoi : M4, convexité par le signe de la dérivée seconde.}
    \item \textbf{D} — Une dérivée seconde positive n'impose pas $f'\leq0$ ; elle rend $f'$ croissante et ne suffit pas a rendre $f$ décroissante. \emph{Renvoi : M4, convexité par le signe de la dérivée seconde.}
  \end{itemize}
  \item \textbf{Q11}
  \begin{itemize}
    \item \textbf{A} — $f'(x)=0$ caractérise un point critique lorsque $f$ est dérivable, pas un changement de convexité. \emph{Renvoi : cours C5, définition.}
    \item \textbf{B} — $f(x)=0$ indique une intersection avec l'axe des abscisses, sans information sur la convexité. \emph{Renvoi : cours C5, définition.}
    \item \textbf{C} — $f''(x)=0$ ne suffit pas : par exemple $f(x)=x^4$ vérifie $f''(0)=0$ sans changer de convexité en 0. \emph{Renvoi : cours C5, point d'inflexion.}
  \end{itemize}
  \item \textbf{Q12}
  \begin{itemize}
    \item \textbf{B} — La position est inversee : une fonction convexe a sa courbe au-dessus de chaque tangente, et non au-dessous. \emph{Renvoi : M4, position de la courbe, de ses tangentes et de ses cordes.}
    \item \textbf{C} — Pour deux points de la courbe d'une fonction convexe, le segment de corde se situe au-dessus de la courbe, pas au-dessous. \emph{Renvoi : M4, position de la courbe, de ses tangentes et de ses cordes.}
    \item \textbf{D} — La première partie est vraie, mais la courbe est au-dessous de ses cordes ; l'affirmation conjointe est donc fausse. \emph{Renvoi : M4, position de la courbe, de ses tangentes et de ses cordes.}
  \end{itemize}
  \item \textbf{Q13}
  \begin{itemize}
    \item \textbf{A} — La condition $f'(0)=0$ indique seulement une tangente horizontale ; elle n'implique pas que la courbe traverse cette tangente. \emph{Renvoi : M4, point d'inflexion et changement de convexité.}
    \item \textbf{C} — L'égalité $f(0)=0$ place le point sur l'axe des abscisses, mais ne donne aucune information sur un changement de convexité. \emph{Renvoi : M4, point d'inflexion et changement de convexité.}
    \item \textbf{D} — La fonction $x^3$ est croissante sur $\mathbb{R}$ ; surtout, la monotonie ne caractérise pas le franchissement d'une tangente. \emph{Renvoi : M4, point d'inflexion et changement de convexité.}
  \end{itemize}
  \item \textbf{Q14}
  \begin{itemize}
    \item \textbf{A} — L'égalité avec la tangente pour tous $a,x$ caracteriserait une fonction affine ; la convexité ne donne qu'une inégalité. \emph{Renvoi : M3, inégalité de convexité par la tangente.}
    \item \textbf{B} — Le sens $\leq$ correspondrait a une fonction concave ; une fonction convexe reste au-dessus de sa tangente en $a$. \emph{Renvoi : M3, inégalité de convexité par la tangente.}
    \item \textbf{D} — La convexité compare $f(x)$ a la tangente $f(a)+f'(a)(x-a)$, pas a la seule valeur $f(a)$ ; aucune croissance n'est supposee. \emph{Renvoi : M3, inégalité de convexité par la tangente.}
  \end{itemize}
  \item \textbf{Q15}
  \begin{itemize}
    \item \textbf{A} — La convexité déterminé le sens de l'approximation : la courbe est au-dessus de toute tangente, donc $1{,}1$ est une valeur par defaut. \emph{Renvoi : M3, sens d'une approximation affine par convexité.}
    \item \textbf{B} — La fonction exponentielle est strictement convexe : sa courbe est au-dessus de sa tangente en $0$, donc $\mathrm{e}^{0{,}1}>1{,}1$. \emph{Renvoi : M3, sens d'une approximation affine par convexité.}
    \item \textbf{C} — La tangente ne coincide avec la courbe qu'au point de contact $x=0$ ; a $x=0{,}1$, la stricte convexité exclut l'égalité. \emph{Renvoi : M3, sens d'une approximation affine par convexité.}
  \end{itemize}
\end{itemize}
\fi
% NEXUS-QCM-TEACHER-ONLY-END
